\section{Problem Statement and Network Infrastructure Architecture}

\subsection{Operational Network Challenges}
The deployment and maintenance of Network Security and Threat Mitigation in complex organizational topologies encounters Unencrypted legacy internal communication protocols, broadcast domain poisoning via ARP spoofing, man-in-the-middle SSL stripping, and credential extraction.. The complete system architecture must maintain high packet delivery reliability, protect data in transit via cryptographic primitives, and deliver deterministic response latency.

\subsection{System Topology and Engineering Pipeline}
Figure~\ref{fig:pipeline} depicts the modular network architecture and diagnostic pipeline designed for this case study. It spans Ingress Filtering, Deep Packet Inspection, Transport Optimization, Routing Convergence, and Telemetry Monitoring.

\begin{figure}[H]
\centering
\resizebox{\textwidth}{!}{\begin{tikzpicture}[
    block/.style={rectangle, draw=black!80, fill=red!8, thick, rounded corners, minimum height=1.1cm, text width=2.5cm, font=\footnotesize},
    arr/.style={-{Stealth[length=2mm]}, thick}
]
\node[block, align=center] (n1) at (0,0) {1. Adversary Host\\Cain / Ettercap ARP Poison};
\node[block, align=center] (n2) at (3.2,0) {2. Unprotected L2 Switch\\Broadcast Flood \& CAM Crash};
\node[block, align=center] (n3) at (6.4,0) {3. Man-in-the-Middle\\Traffic Redirection to Attacker};
\node[block, align=center] (n4) at (9.6,0) {4. Dynamic ARP Inspection\\DAI Binding Verification};
\node[block, align=center] (n5) at (12.8,0) {5. DHCP Snooping Guard\\Untrusted Port Blockade};
\node[block, align=center] (n6) at (12.8,-2.2) {6. IPsec / WireGuard Tunnel\\Encapsulating Security ESP};
\node[block, align=center] (n7) at (9.6,-2.2) {7. TLS 1.3 Strict Policy\\Perfect Forward Secrecy};
\node[block, align=center] (n8) at (6.4,-2.2) {8. 802.1X Port Control\\RADIUS EAP-TLS Auth};
\node[block, align=center] (n9) at (3.2,-2.2) {9. Zero-Trust Microseg.\\Private VLANs (PVLAN)};
\node[block, align=center] (n10) at (0,-2.2) {10. Hardened Network\\Zero Lateral Exploit};
\draw[arr] (n1) -- (n2);
\draw[arr] (n2) -- (n3);
\draw[arr] (n3) -- (n4);
\draw[arr] (n4) -- (n5);
\draw[arr] (n5) -- (n6);
\draw[arr] (n6) -- (n7);
\draw[arr] (n7) -- (n8);
\draw[arr] (n8) -- (n9);
\draw[arr] (n9) -- (n10);
\end{tikzpicture}}
\caption{System Architecture and Engineering Pipeline for Network Security and Threat Mitigation.}
\label{fig:pipeline}
\end{figure}

\subsection{Architectural Component Specifications}
\begin{enumerate}
  \item \textbf{Border Routing \& ACL Layer:} Filters spoofed IP blocks (BCP 38 anti-spoofing) and applies coarse stateless line-rate access lists at the WAN perimeter.
  \item \textbf{Stateful Inspection \& Threat Engine:} Maintains bi-directional connection tracking tables, inspecting Layer 4 flags and protocol headers for anomalies.
  \item \textbf{Deep Packet Inspection (DPI)::} Snort / Suricata engine parses application-layer payloads against rule databases, extracting malicious shellcode patterns.
  \item \textbf{Core Switching \& Routing Fabric:} Multi-layer Gigabit switches enforce 802.1Q VLAN isolation, Spanning Tree (RSTP), and OSPF area routing.
  \item \textbf{Diagnostic Telemetry Sensor:} Continuous NetFlow/sFlow probes and mirrored SPAN ports deliver real-time metrics to network monitoring systems.
\end{enumerate}
